% !TEX root = ../main.tex
\documentclass[../main.tex]{subfiles}

\begin{document}

\intermezzo{On Multimodal Cartography}

\begin{center}
\begin{tikzpicture}
    \fill[CoverGold] (-2,0) rectangle (-1.7,0.08);
    \fill[AccentPurple] (-1.5,0) rectangle (-1.2,0.08);
    \fill[AccentCyan] (-1,0) rectangle (-0.7,0.08);
    \fill[PandaBamboo] (-0.5,0) rectangle (-0.2,0.08);
    \fill[PandaRose] (0,0) rectangle (0.3,0.08);
\end{tikzpicture}
\end{center}

\vspace{1.5cm}

\begin{adjustwidth}{1cm}{1cm}

\noindent\textit{\large Before we descend into methods and data, a pause. What does it mean to make maps of something as fluid as roles? Why insist on multiple modes of seeing? And what ancient wisdom does the cartographer's craft offer the builder of research tools?}

\vspace{1.5cm}

\section*{Cartography: The Writing of Territory}

The word arrives from the Greek: \textit{khártēs} (χάρτης), meaning papyrus or sheet, and \textit{graphein} (γράφειν), to write or draw. \textbf{Cartography} is literally the \textit{writing of sheets}—the inscription of knowledge onto surfaces that can be carried, shared, consulted.

\sidenote{\textit{khártēs} (χάρτης)---papyrus, sheet. \textit{graphein} (γράφειν)---to write, to scratch, to draw.}

But the etymology conceals a deeper truth. The Greeks used \textit{graphein} for both writing and drawing—no distinction between text and image, between description and depiction. A map was always both: a picture \textit{and} a statement, a representation \textit{and} an argument.

This fusion matters for what we attempt. Role constellation cartography is not merely visualization (though it produces images) nor merely analysis (though it generates claims). It is \textit{graphein} in the ancient sense: the inscription of understanding onto surfaces that make the invisible navigable.

\vspace{0.8cm}

\section*{Mappa Mundi: The Napkin of the World}

The Latin \textit{mappa} meant cloth or napkin—something humble, portable, useful. A \textit{mappa mundi} was literally a ``napkin of the world,'' a piece of fabric onto which the known world was sketched for travelers and scholars.

\sidenote{\textit{mappa} (Latin)---cloth, napkin, signal-flag. Hence: something portable, practical, meant to be \textit{used}.}

Consider what this implies. Maps were never meant to be complete representations. They were \textit{tools}—aids to navigation, memory, and decision. The medieval \textit{mappa mundi} placed Jerusalem at the center not from ignorance but from \textit{purpose}: it served pilgrims, not merchants. Different purposes demanded different centers.

Our maps of role constellations inherit this pragmatic lineage. We do not seek the ``true'' map—there is none. We seek \textit{useful} maps: representations that serve navigation through complex innovation ecosystems. The center we choose—the features we emphasize—reflects the journeys we enable.

\vspace{0.8cm}

\section*{The Cartographer's Dilemma}

Every map is a lie that tells the truth. It flattens, simplifies, and selects—yet in doing so, it reveals what was invisible. The London Tube map bears no resemblance to geographic reality, yet it \textit{works}. It serves navigation, not surveying.

\sidenote{\textit{``The map is not the territory.''---Alfred Korzybski.} \textit{``But a good map makes the territory traversable.''---Every cartographer.}}

Harry Beck's 1931 Tube map committed cartographic heresy: he abandoned geography entirely. Distances became meaningless; only \textit{connections} mattered. The map was false to the land but true to the \textit{task}—getting from here to there via the network.

Role constellation maps face the same dilemma. We cannot capture everything. The moment we choose what to represent—connections, positions, claims, flows—we exclude something else. The question is not whether to simplify, but how to simplify \textit{usefully}.

This is where Design Science principles enter. Herbert Simon's insight that design is about ``transforming existing situations into preferred ones'' applies directly: the cartographer transforms invisible complexity into navigable representation. The quality of the transformation is judged not by completeness but by \textit{utility for purpose}.

\vspace{0.8cm}

\section*{Topos and Logos: The Logic of Place}

The Greeks gave us \textit{topos} (τόπος)—place, location, position. From it derives \textit{topology}: the study of spatial properties preserved under continuous deformation. A coffee cup and a donut are topologically identical; what matters is connection, not shape.

\sidenote{\textit{topos} (τόπος)---place, position, topic. \textit{logos} (λόγος)---word, reason, account. Topology: the \textit{logos} of \textit{topos}---the logic of place.}

Role constellations have topology. Organizations occupy \textit{positions} relative to each other—not in physical space but in functional space, relational space, identity space. Some positions are ``close'' (similar functions, overlapping networks); others are ``distant.'' Some organizations bridge otherwise disconnected regions; others cluster at the core.

Network science formalizes this intuition. But the insight predates the formalism. When we map constellations, we are practicing \textit{topology} in the original sense: discerning the logic of positions, the grammar of proximity and distance that shapes what actors can see, reach, and become.

\vspace{0.8cm}

\section*{Why Multiple Maps?}

A single map serves a single purpose. The navigator needs routes; the geologist needs strata; the urban planner needs density. No map serves all purposes.

The \textit{portolan charts} of medieval Mediterranean sailors showed coastlines in exquisite detail but ignored the interior entirely. The \textit{T-O maps} of medieval scholars showed theological order but couldn't guide a ship. Each was ``true'' to its purpose, ``false'' to others.

\sidenote{\textit{Five maps. Five truths. One constellation.}}

Roles, too, have multiple dimensions that no single map captures:

\begin{itemize}[leftmargin=2em, itemsep=0.3em]
    \item What actors \textit{claim} to be (the \textbf{declared} map)
    \item What others \textit{perceive} them to be (the \textbf{attributed} map)
    \item What they \textit{materially} do (the \textbf{enacted} map)
    \item Who they \textit{connect} with (the \textbf{relational} map)
    \item How they \textit{present} themselves (the \textbf{aesthetic} map)
\end{itemize}

A startup's website proclaims ``disruptive innovation.'' News coverage calls them ``promising but unproven.'' Investment data shows modest seed funding. Twitter reveals connections mainly within their own sector. Visual branding mimics established players.

Which is the ``real'' role? All of them. None of them. The role is the \textit{pattern across} these perspectives—the convergences and the telling divergences.

This is why \textbf{multimodal} cartography is not methodological excess but epistemological necessity. \textit{Modus} in Latin meant manner, measure, way. Multiple \textit{modi}—multiple ways of measuring, multiple manners of seeing—are required because the phenomenon itself is multi-aspected. We do not impose multiplicity; we respond to it.

\vspace{0.8cm}

\section*{Convergence and Divergence}

When our five maps align, we gain confidence. An organization claiming intermediary status, recognized as such by media, showing bridging patterns in networks, and channeling resources across boundaries—this is robust evidence. The convergence across \textit{modi} is itself informative: it suggests a role that is not merely performed in one register but \textit{inhabited} across contexts.

\sidenote{\textit{con-vergere}---to incline together. \textit{di-vergere}---to incline apart. Both movements carry signal.}

But divergence is equally informative. When claims don't match evidence, we've found something interesting:

\begin{quote}
\textit{Aspiration gaps}—what actors want to become\\
\textit{Perception gaps}—how observers misread signals\\
\textit{Performance gaps}—where intention exceeds capacity\\
\textit{Strategic gaps}—deliberate misdirection or positioning
\end{quote}

The map that shows only convergence is incomplete. The full cartography includes the tensions—the places where declared and enacted roles pull apart, where relational position contradicts aesthetic positioning. These tensions are not noise; they are \textit{data about dynamics}. They reveal where the constellation is in motion, where roles are contested, where the future is being negotiated.

\vspace{0.8cm}

\section*{On Digital Traces: \textit{Vestigium} and \textit{Signum}}

We map through traces left behind. Website text, news articles, funding records, tweets, images. The Latin offers two relevant words: \textit{vestigium} (footprint, track—something left unintentionally) and \textit{signum} (sign, signal—something made intentionally).

\sidenote{\textit{vestigium}---footprint, trace, track. \textit{signum}---sign, signal, mark. Digital traces are both: unintended residues \textit{and} deliberate signals.}

Digital traces are both. A website is a \textit{signum}—a deliberate signal crafted for audiences. But its metadata, its link structure, its archival history are \textit{vestigia}—footprints left in the course of signaling, often unconsidered. News coverage is someone else's \textit{signum} about the organization—a signal that becomes, for us, a trace of how others perceive.

This dual nature is not a limitation to apologize for. It is the nature of social observation. Interviews, too, are performances. Surveys elicit strategic responses. Ethnography is shaped by the observer's presence. All social data are \textit{signa} from some angle, \textit{vestigia} from another.

Digital traces have their own biases—toward the visible, the English-speaking, the digitally active, the organizationally formalized. We acknowledge these. But they also offer something rare: scale without entirely sacrificing texture. We can read thousands of websites, not dozens. We can trace networks across continents. We can observe dynamics over years through archived \textit{vestigia}.

\vspace{0.8cm}

\section*{Instrumentum: The Tools We Build}

Here Design Science enters explicitly. The Latin \textit{instrumentum} meant tool, equipment, apparatus—but also document, deed, means of proof. An instrument was something built to serve a purpose, and also something that provided evidence.

\sidenote{\textit{instrumentum}---tool, document, means. The \RoleBox and \RoleSim are \textit{instrumenta} in both senses: tools for seeing \textit{and} evidence for claims.}

The \textbf{\RoleBox} platform is an \textit{instrumentum} in this double sense. It is a tool—software that collects, processes, integrates, and visualizes multimodal data about role constellations. But it is also a document of methodological choices: which traces to collect, how to process them, what features to extract, how to integrate across \textit{modi}. The tool \textit{embodies} theoretical commitments; examining the tool reveals the theory.

Design Science Research (DSR) makes this embodiment explicit. Where traditional research treats tools as neutral means to theoretical ends, DSR recognizes that \textit{building the tool is itself a knowledge contribution}. The \RoleBox architecture—its three layers, its modular pipelines, its integration protocols—constitutes a claim about \textit{how} multimodal cartography should be done. The claim is tested not only by whether the maps prove useful but by whether the tool can be replicated, adapted, extended by others.

\vspace{0.5cm}

% Using center instead of figure to avoid kaobook caption recursion
\begin{center}
\begin{tikzpicture}[scale=0.75]
    % DSR cycle adapted for cartography
    \node[draw, circle, fill=Azure!20, minimum size=1.3cm, font=\scriptsize, align=center]
          (prob) at (90:2.5) {Territory\\unknown};
    \node[draw, circle, fill=PandaBamboo!20, minimum size=1.3cm, font=\scriptsize, align=center]
          (obj) at (18:2.5) {Navigation\\needed};
    \node[draw, circle, fill=Marigold!20, minimum size=1.3cm, font=\scriptsize, align=center]
          (des) at (-54:2.5) {Map\\designed};
    \node[draw, circle, fill=AccentPurple!20, minimum size=1.3cm, font=\scriptsize, align=center]
          (demo) at (-126:2.5) {Journey\\attempted};
    \node[draw, circle, fill=AccentCoral!20, minimum size=1.3cm, font=\scriptsize, align=center]
          (eval) at (-198:2.5) {Map\\refined};

    % Arrows
    \draw[->, thick, PandaCharcoal!70] (prob) -- (obj);
    \draw[->, thick, PandaCharcoal!70] (obj) -- (des);
    \draw[->, thick, PandaCharcoal!70] (des) -- (demo);
    \draw[->, thick, PandaCharcoal!70] (demo) -- (eval);
    \draw[->, thick, PandaCharcoal!70] (eval) -- (prob);

    % Center
    \node[font=\scriptsize\bfseries] at (0,0) {Cartographic};
    \node[font=\scriptsize\bfseries] at (0,-0.4) {DSR};
\end{tikzpicture}

\small\textit{The Design Science cycle as cartographic practice: unknown territory prompts need for navigation; need motivates map design; design enables journey; journey reveals map limitations; limitations refine understanding of territory. The cycle continues.}\normalsize
\end{center}

\vspace{0.5cm}

\section*{Simulacrum: The Map That Moves}

But static maps have limits. They show where things \textit{are}, not where they \textit{could be}. They depict the territory as found, not as it might become under different conditions. For this, we need a different kind of representation.

\sidenote{\textit{simulacrum}---likeness, image, imitation. From \textit{simulare}: to make like, to represent, to feign. A simulation is a \textit{simulacrum} that moves through time.}

The Latin \textit{simulacrum} meant image, likeness, or imitation—something made to resemble something else. The English ``simulation'' carries this forward: a \textit{simulacrum} that unfolds dynamically, that responds to perturbation, that generates possible futures rather than depicting the present.

The \textbf{\RoleSim} agent-based model is our dynamic \textit{simulacrum}. Where \RoleBox maps the constellation as observed, \RoleSim asks: \textit{What if?} What if a key intermediary exited? What if policy conditions shifted? What if the fitness landscape tilted toward different role configurations? These counterfactual territories cannot be observed—but they can be \textit{simulated}.

This is not prediction in the forecasting sense. Complex Adaptive Systems resist point prediction; that is their nature. Rather, \RoleSim explores the \textit{space of possibilities}—the adjacent territories that the constellation might traverse, the paths not taken that illuminate why the actual path was taken. It is cartography of the \textit{potential}, complementing cartography of the \textit{actual}.

\vspace{0.5cm}

% Using center instead of figure to avoid kaobook caption recursion
\begin{center}
\begin{tikzpicture}[scale=0.8]
    % RoleBox: Static map
    \begin{scope}[xshift=-4cm]
        \node[font=\small\bfseries, text=Azure] at (0, 2) {\RoleBox};
        \node[font=\scriptsize\itshape, text=Slate] at (0, 1.5) {static cartography};

        % Grid representing mapped territory
        \draw[Azure!30, step=0.4] (-1.2,-0.8) grid (1.2,0.8);

        % Points representing observed positions
        \node[draw, circle, fill=Azure!60, minimum size=0.2cm] at (-0.6, 0.4) {};
        \node[draw, circle, fill=Azure!60, minimum size=0.2cm] at (0.2, 0.2) {};
        \node[draw, circle, fill=Azure!60, minimum size=0.2cm] at (0.6, -0.2) {};
        \node[draw, circle, fill=Azure!60, minimum size=0.2cm] at (-0.2, -0.4) {};

        \node[font=\tiny, text=Slate] at (0, -1.3) {Where things ARE};
    \end{scope}

    % RoleSim: Dynamic simulation
    \begin{scope}[xshift=4cm]
        \node[font=\small\bfseries, text=Marigold] at (0, 2) {\RoleSim};
        \node[font=\scriptsize\itshape, text=Slate] at (0, 1.5) {dynamic simulacrum};

        % Contour lines representing possibility space
        \draw[Marigold!40] (0,0) ellipse (1.2 and 0.8);
        \draw[Marigold!30] (0,0) ellipse (0.8 and 0.5);
        \draw[Marigold!20] (0,0) ellipse (0.4 and 0.25);

        % Trajectory arrows
        \draw[->, Marigold!70, thick] (-0.3, 0.1) -- (0.2, 0.3);
        \draw[->, Marigold!70, thick] (0.2, 0.3) -- (0.5, 0.1);
        \draw[->, Marigold!50, dashed] (0.2, 0.3) -- (0.1, 0.6);
        \draw[->, Marigold!50, dashed] (0.2, 0.3) -- (0.6, 0.4);

        \node[font=\tiny, text=Slate] at (0, -1.3) {Where things COULD BE};
    \end{scope}

    % Connecting arrow
    \draw[<->, thick, PandaCharcoal] (-1.8, 0) -- (1.8, 0);
    \node[font=\tiny, text=PandaCharcoal, align=center] at (0, 0.4) {parameterization};
    \node[font=\tiny, text=PandaCharcoal, align=center] at (0, -0.4) {validation};
\end{tikzpicture}

\small\textit{\RoleBox maps the observed constellation; \RoleSim simulates possible trajectories. The observed parameterizes the possible; the possible contextualizes the observed.}\normalsize
\end{center}

\vspace{0.8cm}

\section*{The Craft of Integration: \textit{Integrare}}

The hardest part is not collecting five types of data. It is \textit{integrating} them into coherent understanding. The Latin \textit{integrare} meant to make whole, to renew, to restore—from \textit{integer}, meaning untouched, whole, complete.

\sidenote{\textit{integrare}---to make whole, restore. \textit{integer}---whole, complete, untouched. Integration is the work of making wholeness from fragments.}

Integration is not mere aggregation. It is not averaging signals or stacking data. It is the work of \textit{making whole}—discerning how fragments relate, where they cohere, where they productively conflict. How do you combine a network position with a visual aesthetic? A funding pattern with a discourse style?

There is no algorithm for this. There are heuristics, frameworks, systematic comparisons. But ultimately, cartography is a \textit{craft}—skilled judgment about what patterns mean and how to represent them. The Latin \textit{ars} (whence ``art'') meant skill, craft, technique—something learned through practice, refined through repetition, never fully reducible to rules.

This dissertation offers tools. It does not offer automation. The human interpreter remains essential—not as a weakness, but as the source of meaning. Computers can detect patterns; humans discern \textit{significance}. The integration of multimodal signals into claims about roles requires both.

\vspace{0.8cm}

\section*{Design Principles as Cartographic Commitments}

The four Meta-Design Requirements introduced in Chapter 1 are, at root, cartographic commitments:

\begin{description}[leftmargin=1em, itemsep=0.4em]
    \item[MDR-1: Theory-Data Reconciliation] \hfill\\
    \textit{The map must speak the language of both terrain and traveler.}\\
    Theoretical concepts (emergence, self-organization) must connect to observable traces; observed patterns must be interpretable through theoretical lenses. Neither theory-blind empiricism nor data-free theorizing suffices.

    \item[MDR-2: Hypothesis-Driven Data Selection] \hfill\\
    \textit{The cartographer chooses what to survey based on the journey's purpose.}\\
    We do not map everything; we map what the fifteen propositions require. Data collection serves theoretical testing, not exhaustive documentation.

    \item[MDR-3: Computational-Human Complementarity] \hfill\\
    \textit{The instrument extends the eye; it does not replace the mind.}\\
    Algorithms detect; humans interpret. Scale requires computation; meaning requires judgment. The \textit{instrumentum} amplifies human capacity without substituting for it.

    \item[MDR-4: Knowledge Accumulation Infrastructure] \hfill\\
    \textit{Good maps enable others to journey—and to make better maps.}\\
    Open tools, documented methods, archived data: these allow replication, extension, critique. The cartographic contribution is not only the maps we make but the map-making capacity we share.
\end{description}

\vspace{0.8cm}

\section*{What Follows}

The chapters ahead will get technical. We will discuss APIs, NLP pipelines, network metrics, image classifiers, agent-based models. We will present tables and graphs and statistics. We will write of recall and precision, of centrality and clustering, of fitness functions and parameter sweeps.

Do not mistake the machinery for the purpose.

\sidenote{\textit{The purpose is seeing---seeing role constellations that would otherwise remain invisible.}}

The purpose is \textit{seeing}—seeing role constellations that would otherwise remain invisible. The machinery serves the seeing. When the machinery obscures the seeing, we have failed.

The ancient cartographers inscribed \textit{hic sunt dracones}—``here be dragons''—on unknown territories. It was honest: an admission of ignorance, a warning to travelers, an invitation to future explorers.

Our maps, too, will have their dragons. Territories where data is sparse. Regions where integration falters. Zones where the traces mislead. We will mark these honestly—not as failures but as frontiers.

\vspace{1cm}

\begin{center}
\textit{Now, to the methods that make the maps.}
\end{center}

\vspace{0.5cm}

\begin{center}
\begin{tikzpicture}
    \fill[CoverGold] (-2,0) rectangle (-1.7,0.08);
    \fill[AccentPurple] (-1.5,0) rectangle (-1.2,0.08);
    \fill[AccentCyan] (-1,0) rectangle (-0.7,0.08);
    \fill[PandaBamboo] (-0.5,0) rectangle (-0.2,0.08);
    \fill[PandaRose] (0,0) rectangle (0.3,0.08);
\end{tikzpicture}
\end{center}

\end{adjustwidth}

\closeintermezzo

\end{document}
